% Section 7: Open problems
\subsection{Open problems}
\label{sec:forecasting_rl_open}

Four problems deserve attention from researchers who already understand
econometric identification. We list them in the order of expected
methodological yield and we sketch the credible attack on each.

The first and most consequential is the construction of action-conditional
time series foundation models. As laid out in Section~\ref{sec:tsfm}, every
existing TSFM at scale, namely Chronos \citep{Ansari2024},
TimesFM \citep{Das2024}, and Lag-Llama \citep{Rasul2023}, conditions the
forecast on observed history but not on the agent's action. The
closest existing work is \citet{Ada2025}, which uses Lag-Llama as a
passive observation-process prior inside an offline-RL loop and combines it
with a conditional diffusion model that handles within-episode dynamics. The
Lag-Llama component does not model the action-conditional transition
$p(y_{t+1} \mid y_{1:t}, a_t)$; the diffusion model does, and only on
stationary offline data. The natural architectural step is a pretraining
objective that explicitly models intervention. The technical questions are
how to define an interventional pretraining objective on observational
data, how to allocate the corpus between observational and interventional
data, and how to validate the resulting forecaster's action-conditional
calibration on held-out interventional benchmarks. The Drain Model of
\citet{Amazon2025} is the closest current production example of an
action-conditional probabilistic simulator, but it is trained on
production-system data with explicit policy choices encoded and is not a
pretrained foundation model.

The second is identification under TSFM-induced counterfactuals. The
structural-causal-model programme of \citet{Buesing2019} and
\citet{Oberst2019} assumes a correctly specified causal graph, namely that
the postulated exogenous noise is independent of past observations
conditional on the state and action. The Lucas critique applies in
its sharpest form when the agent's policy itself enters the data-generating
process, because the SCM's noise marginals depend on the deployed policy
through equilibrium considerations. The credible remedies are imported from
the econometric identification toolkit. Instrumental variables, in the
spirit of \citet{Chernozhukov2018} and the broader
machine-learning-instruments literature, can restore identification of the
action-conditional dynamics under exogeneity restrictions. The most
plausible source of instruments in macroeconomic settings is the
high-frequency monetary policy shock identification of Romer-Romer and the
shadow-rate identification of Wu-Xia. In micro settings, the source is
typically a randomised experiment or a quasi-experimental discontinuity.
The methodological question is how to embed an instrumental variable
restriction into the TSFM pretraining objective, which to our knowledge no
existing work has addressed.

The third is online decision-focused regret in non-stationary economic
environments. \citet{Capitaine2025} provide the first sublinear regret
bounds for online DFL, with the DF-FTPL algorithm attaining static regret
$\widetilde O(m^{3/4} n^{1/2} T^{-1/4})$ and the DF-OGD algorithm attaining
dynamic regret $\widetilde O((1 + P_T)^{1/4} T^{-1/4} n^{-1/2})$, both
under a margin condition on the optimal vertex of the decision polytope.
The bounds depend on the path variation $P_T$ of the comparator, which is
the right quantity for non-stationary environments and is exactly the
quantity that the macroeconomic shock literature spends most of its energy
estimating. The open problem is the application to economic data. The
synthetic knapsack experiment of \citet{Capitaine2025} is far from a
demand-forecasting or asset-pricing benchmark on real data, and the gap
between the formal regret rate and the empirical performance on noisy
macroeconomic series remains uncharted. The natural project is to apply
DF-OGD to the contextual newsvendor benchmark of \citet{BanRudin2019}, to
the Cao-Xu microgrid task of \citet{CaoXu2025}, or to the macro reaction
function of \citet{HinterlangTaenzer2024}, and to measure both the
realised regret and the path variation $P_T$ on the historical data.

The fourth is form-aware pretraining objectives so that TSFMs produce
joint trajectory ensembles rather than the per-step marginal point or
parametric forecasts that the form-and-function critique of
\citet{PerezDiaz2025} identifies as the current operational limitation.
Their Propositions 1-3 establish that path-dependent operational queries,
including threshold-crossing inventory rules, value-at-risk over a
multi-period horizon, and run-length distributions for predictive
maintenance, are not answerable from marginally-calibrated forecasts
without imposing a copula or another dependence structure that the data
may not support. The credible remedy is a pretraining objective that
elicits joint calibration directly, namely a strictly proper scoring rule
on the joint trajectory distribution. The Energy Score is the natural
candidate, but its computational cost scales with the trajectory length
squared and the joint-Energy-Score-trained variants of Chronos or TimesFM
do not yet exist at scale. A second route is to convert any marginally
trained TSFM into a trajectory ensemble by autoregressive sampling and to
fine-tune the resulting ensemble on a small dataset of path-dependent
queries, in the spirit of decision-focused fine-tuning. The third and
most ambitious route is to redesign the TSFM architecture so that the
joint-distribution head is the default output, with the marginal heads
derived as projections. The diffusion-trajectory approach of
\citet{Janner2022diffuser} and \citet{Ajay2023} is the closest existing
template, although applied to RL trajectories rather than to univariate
time series.

The four open problems are linked by a single methodological observation.
The forecasting side of the chapter has matured into a stack of pretrained
models with impressive marginal accuracy, and the decision side has matured
into a stack of differentiable decision-focused training procedures with
quantitative finite-sample guarantees. The interface between the two, namely
the action-conditional and identification-aware foundation forecaster
trained under a joint-calibration objective inside an online decision-focused
loop, does not yet exist. An economist who understands the Lucas critique,
the proper-scoring-rule literature, and the instrumental-variable toolkit is
in an unusually good position to build it. The simulation in
Section~\ref{sec:forecasting_rl_sim} is a deliberate stripped-down
demonstration of the same principle on the simplest econometric substrate.
Two forecasts that look identical on the MSE metric differ by a factor of
nearly five on realised regret once the noise model is misspecified, and
the corrective is a training objective that knows about the decision. The
generalisation of that recipe to action-conditional pretraining, to
identified counterfactual forecasting, to online non-stationary regret, and
to joint trajectory ensembles is the chapter's research agenda for the next
five years.
